\begin{examples}[7.5.3]
\label{IV.2.s7.5.3-fr}
Consideremos las siguientes propiedades $\boldsymbol R(C)$, (siendo $C$ un \emph{anillo
local noetheriano}):
\begin{enumerate}
\item[(i)] (También denotada por $(\mathrm i_n)$)
$\operatorname{coprof}(C)\leq n$.
\oldpage[IV-2]{206}
\item[(ii)] (También denotada por $(\mathrm{ii}_n)$) $C$
satisface $(\mathrm S_n)$.
\item[(iii)] $C$ es un anillo de Cohen-Macaulay.
\item[(iv)] $C$ es un anillo regular.
\item[(v)] (También denotada por $(\mathrm v_n)$) $C$ es íntegro,
integralmente cerrado y satisface $(\mathrm R_n)$.
\item[(vi)] $C$ es íntegro e integralmente cerrado.
\item[(vii)] $C$ es íntegro.
\item[(viii)] $C$ es reducido.
\end{enumerate}
Todas estas propiedades satisfacen $(\mathrm R_{\mathrm{III}})$,
como se ha visto en (7.3.19, (i)). También satisfacen
$(\mathrm R_{\mathrm{IV}})$: para (i), esto resulta de
(\textbf 0, 16.4.10, (ii)); para (ii) y (iii), resulta de (5.12.4) y del hecho
de que un anillo local noetheriano completo es catenario (5.6.4);
para (iv), es un caso particular de (\textbf 0, 17.1.8); para (vii), se deduce
de (3.4.5), y para (viii), de (3.4.6); para (vi), resulta de (5.12.7)
y de que un anillo local noetheriano completo es catenario (5.6.4);
por último, para (v), resulta de lo anterior y de (5.12.5).

Se puede, por tanto, aplicar el teorema (7.5.1) cuando $\boldsymbol R$
es cualquiera de las propiedades anteriores. Por otra parte, ya
se ha observado (7.3.11) que las propiedades (i) a (viii) satisfacen
$(\mathrm R_{\mathrm{II}})$ (salvo para (vii), donde la verificación de
$(\mathrm R_{\mathrm{II}})$ es consecuencia de (2.1.14)). En cuanto a
$(\mathrm R_{\mathrm V})$, se observará que, si se toma
$\boldsymbol R'=\boldsymbol R$, la condición
$(\mathrm R_{\mathrm V})$ se reduce a la condición
$(\mathrm R'_{\mathrm I})$ de (7.3.11), y ya se ha señalado que las
propiedades (ii) a (vi), así como (viii), satisfacen
$(\mathrm R'_{\mathrm I})$ (7.3.11). Para la propiedad (i), se puede
tomar como $\boldsymbol R'$ la propiedad de ser un anillo de
Cohen-Macaulay, en virtud de (6.3.2). Por último, para la propiedad
(vii), la condición $(\mathrm R'_{\mathrm I})$ no se satisface
(ni tampoco $(\mathrm R_{\mathrm I})$), como muestra el ejemplo de
(6.15.11, (ii)) o el de (6.5.5, (ii)). En cambio,
$(\mathrm R_{\mathrm V})$ sí se cumple si se toma como
$\boldsymbol R'$ la propiedad de ser \emph{regular}: basta, en efecto,
aplicar el lema (7.5.1.1) tomando como $\boldsymbol{C}$ la categoría de
todos los anillos locales noetherianos y como $\boldsymbol R(C)$ la
propiedad de ser íntegro; esto es posible en virtud de (3.4.5).

Se ve así, en particular, que, cuando $\boldsymbol R$ es cualquiera
de las propiedades (i) a (viii), la conclusión del corolario (7.5.2)
es aplicable si se supone que \emph{las fibras formales de} $A$ \emph{son
geométricamente regulares} (cf. (7.8.2)).
\end{examples}
